\section{Kết quả truy vấn và vai trò của đồ thị}
\label{sec:retrieval}

\subsection{Văn bản $\to$ ảnh trên Visual Genome (RQ1)}
Bảng~\ref{tab:main} là kết quả chính trên tập test. Mọi mô hình có tham số học được đều vượt CLIP thuần rõ rệt: \mrr{} tăng khoảng 0,02 (khoảng 18--19\% tương đối) và R@10 tăng khoảng 0,045. Mô hình chính đạt \mrr{} cao nhất, 0,1284 \pmm{} 0,0006; mô hình có thêm cạnh đồng xuất hiện đạt R@1 và R@5 cao nhất. Hình~\ref{fig:mrr} cho thấy khoảng cách giữa CLIP và các mô hình có học lớn hơn nhiều so với khoảng cách giữa các mô hình có học với nhau.

\begin{table}[t]
\centering
\caption{Văn bản $\to$ ảnh trên test Visual Genome (5.000 ảnh, 25.000 truy vấn); trung bình \pmm{} độ lệch chuẩn qua 3 seed}
\label{tab:main}
\scriptsize
\setlength{\tabcolsep}{1.5pt}
\begin{tabular}{@{}l c c c c@{}}
\toprule
\textbf{Mô hình} & \textbf{R@1} & \textbf{R@5} & \textbf{R@10} & \textbf{\mrr} \\
\midrule
CLIP thuần & 0,0584 & 0,1422 & 0,2018 & 0,1076 \\
Không đồ thị & 0,0689\pmm0,0006 & 0,1741\pmm0,0056 & 0,2457\pmm0,0043 & 0,1266\pmm0,0006 \\
Đồ thị xáo & 0,0686\pmm0,0002 & 0,1759\pmm0,0004 & 0,2489\pmm0,0004 & 0,1272\pmm0,0003 \\
+ Đồng xuất hiện & \textbf{0,0709}\pmm0,0009 & \textbf{0,1792}\pmm0,0012 & 0,2469\pmm0,0006 & 0,1280\pmm0,0006 \\
\textbf{CLIP + Graph} & 0,0705\pmm0,0014 & 0,1765\pmm0,0037 & \textbf{0,2492}\pmm0,0015 & \textbf{0,1284}\pmm0,0006 \\
\bottomrule
\end{tabular}
\end{table}

\begin{figure}[t]
\centering
\begin{tikzpicture}
\begin{axis}[
  ybar, width=\linewidth, height=4.4cm, bar width=13pt,
  ymin=0.10, ymax=0.132, ylabel={\mrr{} (test)}, ylabel style={font=\footnotesize},
  symbolic x coords={CLIP,NoG,Shuf,Cooc,Graph},
  xtick=data, x tick label style={font=\scriptsize,align=center,text width=1.3cm},
  xticklabels={CLIP,Không đồ thị,Đồ thị xáo,Đồng xuất hiện,CLIP + Graph},
  y tick label style={font=\scriptsize,/pgf/number format/fixed,/pgf/number format/precision=3,/pgf/number format/use comma},
  nodes near coords, nodes near coords style={font=\tiny,/pgf/number format/fixed,/pgf/number format/fixed zerofill,/pgf/number format/precision=4,/pgf/number format/use comma},
  every node near coord/.append style={yshift=3pt},
  enlarge x limits=0.13, axis lines*=left, ymajorgrids, grid style={gray!25}]
\addplot[fill=blue!45,draw=blue!70!black,error bars/.cd,y dir=both,y explicit]
 coordinates {(CLIP,0.1076)+-(0,0) (NoG,0.1266)+-(0,0.0006) (Shuf,0.1272)+-(0,0.0003) (Cooc,0.1280)+-(0,0.0006) (Graph,0.1284)+-(0,0.0006)};
\end{axis}
\end{tikzpicture}
\caption{\mrr{} trên test Visual Genome; thanh lỗi là độ lệch chuẩn qua 3 seed. Trục tung bắt đầu từ 0,10 để thấy rõ chênh lệch nhỏ.}
\label{fig:mrr}
\end{figure}

Nhóm chạy lại năm thí nghiệm theo giao thức thứ hai để có thêm nDCG@10 và chiều ảnh $\to$ văn bản, với điều kiện cấu hình giữ nguyên, quy tắc được commit trước, test chạy một lần cho mỗi seed và không sửa dòng kết quả cũ. R@K và MRR của lần chạy lại khớp bảng chính với sai khác không quá 0,0001. Ở chiều ảnh $\to$ văn bản, mô hình không được huấn luyện thêm: truy vấn là mỗi ảnh, kho là 25.000 câu, câu đúng là các câu của chính ảnh và mọi câu giống hệt chúng; $\alpha$ của chiều này được chọn riêng trên val. Bảng~\ref{tab:eval2} cho thấy cùng một xu hướng ở cả hai chiều: các mô hình có học vượt CLIP với khoảng cách lớn hơn nhiều lần độ lệch chuẩn. Giữa ba mô hình có đồ thị (thật, xáo, đồng xuất hiện), chênh lệch xấp xỉ độ lệch chuẩn, nên nhóm không kết luận mô hình nào trong ba mô hình này vượt trội. Mô hình không đồ thị thấp hơn ba mô hình có đồ thị khoảng 0,003 \mrr{} ở chiều ảnh $\to$ văn bản (gấp khoảng 5 lần độ lệch chuẩn), nhưng bản xáo cũng đạt mức đó, nên phần hơn này chưa thể quy cho cấu trúc đồ thị thật.

\begin{table}[t]
\centering
\caption{nDCG@10 (văn bản $\to$ ảnh) và kết quả ảnh $\to$ văn bản trên test (3 seed)}
\label{tab:eval2}
\scriptsize
\setlength{\tabcolsep}{2pt}
\begin{tabular}{@{}l c c c c@{}}
\toprule
& \textbf{T$\to$I} & \multicolumn{3}{c}{\textbf{Ảnh $\to$ văn bản}} \\
\cmidrule(lr){2-2}\cmidrule(l){3-5}
\textbf{Mô hình} & \textbf{nDCG@10} & \textbf{R@1} & \textbf{R@10} & \textbf{\mrr} \\
\midrule
CLIP thuần & 0,1169 & 0,1280 & 0,3630 & 0,2068 \\
Không đồ thị & 0,1387\pmm0,0006 & 0,1484\pmm0,0015 & 0,3895\pmm0,0013 & 0,2281\pmm0,0006 \\
Đồ thị xáo & 0,1390\pmm0,0003 & 0,1519\pmm0,0009 & 0,3957\pmm0,0012 & 0,2312\pmm0,0010 \\
+ Đồng xuất hiện & 0,1383\pmm0,0006 & \textbf{0,1522}\pmm0,0008 & 0,3957\pmm0,0016 & \textbf{0,2316}\pmm0,0008 \\
CLIP + Graph & \textbf{0,1391}\pmm0,0007 & 0,1515\pmm0,0008 & \textbf{0,3961}\pmm0,0024 & 0,2315\pmm0,0007 \\
\bottomrule
\end{tabular}
\end{table}

\subsection{Đồ thị đóng góp bao nhiêu? (RQ2)}
Bảng~\ref{tab:main} cho phép tách mức tăng \mrr{} +0,0208 của mô hình chính so với CLIP thành hai phần. Phần do \emph{phép chiếu học được} (CLIP $\to$ không đồ thị) là +0,0190, khoảng 91\%. Phần do \emph{đồ thị} (không đồ thị $\to$ CLIP + Graph) là +0,0018, khoảng 9\%. Cần lưu ý rằng đối chứng không đồ thị dùng learning rate $10^{-3}$, còn mô hình chính dùng $10^{-4}$ (mỗi mô hình lấy giá trị tốt nhất trên val), nên phép tách này chỉ là ước lượng thô. Trên val, đối chứng không đồ thị với cùng learning rate $10^{-4}$ đạt 0,1246 \pmm{} 0,0003, thấp hơn bản dùng $10^{-3}$ (0,1262 \pmm{} 0,0017).

\begin{table}[t]
\centering
\caption{\mrr{} test ghép cặp theo seed giữa đồ thị thật, đồ thị xáo và mô hình không đồ thị}
\label{tab:paired}
\scriptsize
\setlength{\tabcolsep}{2pt}
\begin{tabular}{@{}c c c c c c@{}}
\toprule
\textbf{Seed} & \textbf{Thật} & \textbf{Xáo} & \textbf{Không ĐT} & \textbf{Thật -- xáo} & \textbf{Xáo -- không ĐT} \\
\midrule
0 & 0,1279 & 0,12681 & 0,1259 & +0,00109 & +0,00091 \\
1 & 0,1283 & 0,12735 & 0,1269 & +0,00095 & +0,00045 \\
2 & 0,1291 & 0,12736 & 0,1271 & +0,00174 & +0,00026 \\
\midrule
TB & 0,1284 & 0,12718 & 0,1266 & \textbf{+0,00126}\pmm0,00042 & +0,00054\pmm0,00034 \\
\bottomrule
\end{tabular}
\end{table}

Ablation xáo cạnh (Bảng~\ref{tab:paired}) cho thấy đồ thị thật hơn đồ thị xáo ở cả 3 seed, trung bình +0,00126 \pmm{} 0,00042 (t ghép cặp bằng 5,16 với 2 bậc tự do). Kết quả này phù hợp với giả thuyết rằng một phần nhỏ của mức cải thiện đến từ quan hệ ảnh--khái niệm thật. Tuy nhiên nhóm không coi đây là bằng chứng: chỉ có 3 seed và một bản xáo, hiệu ứng tuyệt đối chỉ khoảng 1\% \mrr, và giá trị t ở đây chỉ phản ánh dao động theo seed huấn luyện. Xét từng chỉ số, chỉ R@1 có mức chênh (+0,0019) lớn hơn độ lệch chuẩn của mô hình chính. Có thêm hai dấu hiệu phụ cùng chiều: trên đồ thị xáo, $\alpha$ được chọn cao hơn (0,1/0,2/0,1 so với 0/0/0,1) và epoch tốt nhất đến muộn hơn (7/9/8 so với 6/8/6). Điều này gợi ý rằng khi cạnh mang ít thông tin, mô hình có thể dựa vào CLIP nhiều hơn.

Có một điểm bất thường nhóm chưa giải thích được. Trên val, chênh lệch giữa đồ thị thật và đồ thị xáo là +0,0040, nhưng trên test chỉ còn +0,0013: bản xáo tăng +0,0031 từ val lên test, trong khi mô hình chính chỉ tăng +0,0003. Tỉ lệ cạnh xáo trùng với cạnh gốc gần như nhau ở val và test (1,78\% và 1,89\%), nên không giải thích được hiện tượng này. Vì chỉ có một bản xáo, nhóm không tách được phần dao động do chính cách xáo.

Thêm quan hệ khái niệm--khái niệm không cải thiện kết quả trên toàn tập: \mrr{} 0,1280 \pmm{} 0,0006, so với 0,1284 \pmm{} 0,0006 của mô hình một lớp. Nhóm báo cáo kết quả âm này đúng theo quy tắc đã chốt trước.

\subsection{Đồ thị giúp ở nhóm truy vấn nào?}
Truy vấn test được chia theo hai cách độc lập, với luật được chốt trên val trước khi chạy test. Theo \emph{độ đặc trưng}, truy vấn \code{specific} có đúng một ảnh đúng (88,7\%) và \code{generic} có từ hai ảnh đúng trở lên (11,3\%). Theo \emph{quan hệ}, truy vấn thuộc nhóm \code{action} nếu câu chứa vị từ hành động (8,9\%), \code{spatial} nếu chứa vị từ không gian (39,7\%), còn lại là \code{none} (51,4\%).

\begin{table}[t]
\centering
\caption{\mrr{} văn bản $\to$ ảnh trên test theo nhóm truy vấn (3 seed); $n$ là số truy vấn}
\label{tab:groups}
\scriptsize
\setlength{\tabcolsep}{1.5pt}
\begin{tabular}{@{}l r c c c c@{}}
\toprule
\textbf{Nhóm} & $n$ & \textbf{CLIP} & \textbf{Không ĐT} & \textbf{Xáo} & \textbf{CLIP+Graph} \\
\midrule
generic & 2.820 & 0,0681 & 0,0914\pmm0,0092 & 0,0951\pmm0,0018 & \textbf{0,1082}\pmm0,0040 \\
specific & 22.180 & 0,1126 & 0,1311\pmm0,0007 & \textbf{0,1313}\pmm0,0005 & 0,1309\pmm0,0011 \\
\midrule
action & 2.217 & 0,1223 & 0,1446\pmm0,0004 & \textbf{0,1468}\pmm0,0003 & 0,1455\pmm0,0012 \\
spatial & 9.934 & 0,1158 & 0,1396\pmm0,0020 & 0,1410\pmm0,0007 & \textbf{0,1444}\pmm0,0008 \\
none & 12.849 & 0,0987 & 0,1134\pmm0,0011 & 0,1131\pmm0,0006 & 0,1130\pmm0,0012 \\
\bottomrule
\end{tabular}
\end{table}

Bảng~\ref{tab:groups} cho thấy đồ thị phát huy tác dụng ở nhóm \code{generic}. Ở nhóm này, đồ thị thật tách rõ khỏi cả hai đối chứng: \mrr{} đạt 0,1082, so với 0,0951 của đồ thị xáo và 0,0914 của mô hình không đồ thị; R@1 tăng khoảng 50\% (0,042 so với 0,027--0,028). Đây là những câu ngắn và chung chung như ``a man'' hay ``white clouds'', có nhiều ảnh đúng. Một cách hiểu hợp lý là vectơ trung bình của 16 khái niệm giúp đẩy lên những ảnh thực sự chứa khái niệm được nhắc tới, thay vì những ảnh chỉ giống về bề mặt. Bảng~\ref{tab:groups-r1} cho thấy R@1 ở hai nhóm \code{generic} và \code{spatial}.

\begin{table}[t]
\centering
\caption{R@1 văn bản $\to$ ảnh trên test cho nhóm \code{generic} và \code{spatial}}
\label{tab:groups-r1}
\scriptsize
\setlength{\tabcolsep}{1.5pt}
\begin{tabular}{@{}l c c c c c@{}}
\toprule
\textbf{Nhóm} & \textbf{CLIP} & \textbf{Không ĐT} & \textbf{Xáo} & \textbf{CLIP+Graph} & \textbf{+ Đồng XH} \\
\midrule
generic & 0,0238 & 0,0280 & 0,0274 & 0,0424 & \textbf{0,0457} \\
spatial & 0,0642 & 0,0750 & 0,0732 & 0,0777 & \textbf{0,0790} \\
\bottomrule
\end{tabular}
\end{table}

Ở nhóm \code{spatial}, đồ thị thật cũng cao hơn hai đối chứng (0,1444 so với 0,1410 và 0,1396) với mức chênh lớn hơn độ lệch chuẩn. Ở các nhóm \code{specific}, \code{action} và \code{none}, các mô hình có học gần như ngang nhau. Phân tích này dựa trên một lần chia đã chốt và chưa có kiểm định theo từng truy vấn (paired bootstrap).

\subsection{Độ nhạy theo trọng số $\alpha$}
Hình~\ref{fig:alpha} vẽ \mrr{} val theo $\alpha$, lấy trung bình 3 seed. Giá trị tốt nhất nằm gần $\alpha = 0$, tức điểm cuối gần như chỉ còn nhánh đồ thị. Điều này không có nghĩa là CLIP bị bỏ qua, vì vectơ CLIP của ảnh vẫn nằm trong $h_i$ qua $W_{\text{self}}$. Đường cong gần như phẳng trong đoạn $[0; 0,2]$ và chỉ giảm rõ khi $\alpha$ tiến về 1. Với đồ thị xáo, đường cong thấp hơn ở mọi $\alpha < 1$. Bảng~\ref{tab:alpha} liệt kê giá trị của từng seed.

\begin{table}[t]
\centering
\caption{\mrr{} val theo $\alpha$ của từng seed (trích một số giá trị trong lưới); in đậm là giá trị được chọn}
\label{tab:alpha}
\scriptsize
\setlength{\tabcolsep}{2.5pt}
\begin{tabular}{@{}c c c c c c c@{}}
\toprule
$\alpha$ & Thật s0 & Thật s1 & Thật s2 & Xáo s0 & Xáo s1 & Xáo s2 \\
\midrule
0,0 & \textbf{0,1286} & \textbf{0,1285} & 0,1258 & 0,1238 & 0,1240 & 0,1237 \\
0,1 & 0,1266 & 0,1266 & \textbf{0,1272} & \textbf{0,1240} & 0,1240 & \textbf{0,1241} \\
0,2 & 0,1267 & 0,1264 & 0,1258 & 0,1240 & \textbf{0,1240} & 0,1236 \\
0,5 & 0,1228 & 0,1236 & 0,1233 & 0,1219 & 0,1229 & 0,1223 \\
1,0 & 0,1061 & 0,1061 & 0,1061 & 0,1061 & 0,1061 & 0,1061 \\
\bottomrule
\end{tabular}
\end{table}

\begin{figure}[t]
\centering
\begin{tikzpicture}
\begin{axis}[width=0.95\linewidth,height=4.2cm,xlabel={$\alpha$},ylabel={\mrr{} val},
  label style={font=\footnotesize}, tick label style={font=\scriptsize},
  xmin=-0.03,xmax=1.03,ymin=0.104,ymax=0.13,legend pos=south west,legend style={font=\scriptsize},
  y tick label style={/pgf/number format/fixed,/pgf/number format/precision=3,/pgf/number format/use comma},
  x tick label style={/pgf/number format/use comma},
  ymajorgrids,grid style={gray!25}]
\addplot[thick,blue!70!black,mark=*] coordinates {(0,0.12763) (0.1,0.12680) (0.2,0.12630) (0.5,0.12323) (1,0.1061)};
\addlegendentry{Đồ thị thật}
\addplot[thick,red!70!black,mark=square*,dashed] coordinates {(0,0.12383) (0.1,0.12403) (0.2,0.12387) (0.5,0.12237) (1,0.1061)};
\addlegendentry{Đồ thị xáo}
\end{axis}
\end{tikzpicture}
\caption{\mrr{} val theo $\alpha$ (trung bình 3 seed); $\alpha = 1$ là CLIP thuần.}
\label{fig:alpha}
\end{figure}

\subsection{Tập công khai thứ hai: Flickr30k}
\label{sec:flickr}
Bốn thí nghiệm của cấp độ 1 (CLIP thuần, không đồ thị, CLIP + Graph, đồ thị xáo) được chạy lại trên Flickr30k. Cấu hình là bản sao cấu hình Visual Genome, chỉ đổi tên thí nghiệm và đường dẫn dữ liệu; một bài kiểm thử tự động xác nhận điều này. Không siêu tham số nào được chỉnh lại. Epoch và $\alpha$ vẫn được chọn trên val, và tập test được chạy đúng một lần cho mỗi seed.

\begin{table}[t]
\centering
\caption{Văn bản $\to$ ảnh trên test Flickr30k (5.000 ảnh, 25.000 truy vấn, 3 seed)}
\label{tab:flickr}
\scriptsize
\setlength{\tabcolsep}{2pt}
\begin{tabular}{@{}l c c c c@{}}
\toprule
\textbf{Mô hình} & \textbf{R@1} & \textbf{R@10} & \textbf{\mrr} & $\alpha$ \\
\midrule
CLIP thuần & 0,4901 & 0,8250 & 0,6045 & 1 \\
Không đồ thị & \textbf{0,5332}\pmm0,0006 & \textbf{0,8556}\pmm0,0011 & \textbf{0,6447}\pmm0,0003 & 0,6 \\
CLIP + Graph & 0,5309\pmm0,0036 & 0,8553\pmm0,0026 & 0,6432\pmm0,0028 & 0,5 \\
Đồ thị xáo & 0,5276\pmm0,0014 & 0,8526\pmm0,0009 & 0,6399\pmm0,0015 & 0,4--0,5 \\
\bottomrule
\end{tabular}
\end{table}

Bảng~\ref{tab:flickr} lặp lại kết luận chính của Visual Genome: gần như toàn bộ mức tăng so với CLIP đến từ phép chiếu học được. \mrr{} test tăng từ 0,6045 lên 0,6447 (+0,040, khoảng 6,6\% tương đối). So với Visual Genome (+0,019, khoảng 18\%), mức tăng tuyệt đối lớn hơn nhưng mức tăng tương đối nhỏ hơn, vì CLIP vốn đã rất mạnh trên chú thích Flickr30k.

Khác với Visual Genome, trên Flickr30k đồ thị không mang lại thêm gì. Ghép cặp theo seed, CLIP + Graph kém mô hình không đồ thị $-0,0015$ \pmm{} 0,0028 trên test (seed 0 và 1 gần như bằng nhau, seed 2 kém 0,0047 do dừng sớm ở epoch 2). Đồ thị thật vẫn hơn đồ thị xáo ở cả 3 seed (+0,0033 \pmm{} 0,0013), cùng chiều với Visual Genome, nhưng bản xáo lại kém cả mô hình không đồ thị ($-0,0048$ \pmm{} 0,0016). Như vậy trên Flickr30k, cạnh nhiễu gây hại, còn cạnh thật chỉ vừa đủ bù lại phần hại đó. Giá trị $\alpha$ được chọn là 0,4--0,6, nghĩa là điểm cuối chia gần đều giữa CLIP và nhánh học được. Hai hạn chế của thí nghiệm này là từ vựng khái niệm lấy từ Visual Genome nên thiếu nhiều hoạt động của con người vốn hay gặp trong chú thích Flickr30k, và tập train Flickr30k nhỏ hơn Visual Genome khoảng 4,7 lần.

Bảng~\ref{tab:flickr-val} cho kết quả trên val, và Bảng~\ref{tab:flickr-paired} cho hiệu ghép cặp theo từng seed.

\begin{table}[t]
\centering
\caption{Flickr30k trên val (tập dùng để chọn epoch và $\alpha$)}
\label{tab:flickr-val}
\scriptsize
\setlength{\tabcolsep}{2pt}
\begin{tabular}{@{}l c c c c@{}}
\toprule
\textbf{Mô hình} & \textbf{R@1} & \textbf{R@5} & \textbf{R@10} & \textbf{\mrr} \\
\midrule
CLIP thuần & 0,4891 & 0,7368 & 0,8168 & 0,6016 \\
Không đồ thị & 0,5288\pmm0,0003 & 0,7785\pmm0,0010 & 0,8519\pmm0,0008 & 0,6403\pmm0,0000 \\
CLIP + Graph & 0,5298\pmm0,0026 & 0,7782\pmm0,0038 & 0,8517\pmm0,0028 & 0,6408\pmm0,0026 \\
Đồ thị xáo & 0,5269\pmm0,0002 & 0,7753\pmm0,0013 & 0,8495\pmm0,0012 & 0,6383\pmm0,0007 \\
\bottomrule
\end{tabular}
\end{table}

\begin{table}[t]
\centering
\caption{\mrr{} Flickr30k ghép cặp theo seed}
\label{tab:flickr-paired}
\scriptsize
\setlength{\tabcolsep}{2pt}
\begin{tabular}{@{}l l c c c c@{}}
\toprule
\textbf{Tập} & \textbf{Hiệu} & \textbf{Seed 0} & \textbf{Seed 1} & \textbf{Seed 2} & \textbf{TB \pmm{} std} \\
\midrule
Test & Thật -- không ĐT & +0,00032 & $-0,00006$ & $-0,00474$ & $-0,00149$\pmm0,00282 \\
Test & Thật -- xáo & +0,00444 & +0,00354 & +0,00184 & +0,00327\pmm0,00132 \\
Test & Xáo -- không ĐT & $-0,00412$ & $-0,00360$ & $-0,00658$ & $-0,00477$\pmm0,00159 \\
Val & Thật -- không ĐT & +0,00222 & +0,00191 & $-0,00253$ & +0,00053\pmm0,00266 \\
Val & Thật -- xáo & +0,00400 & +0,00328 & +0,00023 & +0,00250\pmm0,00200 \\
Val & Xáo -- không ĐT & $-0,00178$ & $-0,00137$ & $-0,00276$ & $-0,00197$\pmm0,00072 \\
\bottomrule
\end{tabular}
\end{table}

\subsection{Phương pháp năm 2024: G-Retriever (PCST)}
\label{sec:gret}
Để chọn phương pháp mới, nhóm chấm 9 ứng viên (NativE, AdaMF-MAT, Mixed-Curvature MMKGC, ULTRA, CORA, Structure-CLIP, GOAt, G-RAG và G-Retriever) theo mức độ phù hợp với bài toán truy vấn ảnh, khả năng chạy trên cùng pipeline và tình trạng mã nguồn. G-Retriever đạt điểm cao nhất (33/40).

\textbf{Giai đoạn 1.} Cả 3 cấu hình (bản chính, bản không dựng cây, bản xáo cạnh) được chạy trên val với 3 seed. Ở cả 9 lượt, mọi tổ hợp $(k, C)$ đều chọn $\alpha = 1$, tức trùng CLIP thuần (\mrr{} val 0,1061). Nhóm chạy thêm chẩn đoán trên 3.000 câu val đầu tiên để phân biệt lỗi cài đặt với kết quả thật. Không phát hiện lỗi: mọi ảnh val đều có ít nhất một khái niệm và 29--48\% ứng viên có điểm đồ thị dương. Nguyên nhân nằm ở chỗ chênh lệch điểm CLIP giữa hai ứng viên liền kề có trung vị chỉ 0,00037, trong khi điểm ``tỉ lệ chạm đồ thị con'' quá thô để sắp lại những ứng viên sát nhau như vậy (Bảng~\ref{tab:gret-diag}).

\begin{table}[t]
\centering
\caption{Chẩn đoán G-Retriever giai đoạn 1 trên 3.000 câu val đầu tiên (\mrr{} / R@10). $|S_q|$: số khái niệm trong đồ thị con}
\label{tab:gret-diag}
\scriptsize
\setlength{\tabcolsep}{3pt}
\begin{tabular}{@{}c c c@{}}
\toprule
$\alpha$ & $k=3$, $C=0,5$ ($|S_q|\approx 3$) & $k=10$, $C=0,25$ ($|S_q|\approx 12$) \\
\midrule
1,00 (CLIP) & 0,0947 / 0,1697 & 0,0947 / 0,1697 \\
0,99 & 0,0948 / 0,1707 & 0,0945 / 0,1703 \\
0,98 & 0,0949 / 0,1717 & 0,0941 / 0,1710 \\
\bottomrule
\end{tabular}
\end{table}

\textbf{Giai đoạn 2.} Giai đoạn này thay điểm tỉ lệ bằng GraphSAGE trên đồ thị con và được chạy trên cả hai tập với cùng cấu hình. Đối chứng là mô hình không đồ thị có \emph{cùng} cấu hình huấn luyện (lr $10^{-4}$), để mọi khác biệt chỉ đến từ đồ thị con. Bảng~\ref{tab:gret2} cho kết quả trên test, Bảng~\ref{tab:gret2-val} trên val (tập dùng để chọn epoch và $\alpha$), Bảng~\ref{tab:gret2-paired} cho hiệu ghép cặp theo seed.

\begin{table}[t]
\centering
\caption{G-Retriever giai đoạn 2 trên test (5.000 ảnh, 25.000 truy vấn mỗi tập, 3 seed). Đối chứng: không đồ thị, cùng cấu hình huấn luyện (lr $10^{-4}$). Không so số tuyệt đối giữa hai tập}
\label{tab:gret2}
\scriptsize
\setlength{\tabcolsep}{1.5pt}
\resizebox{\columnwidth}{!}{\begin{tabular}{@{}l c c c c c@{}}
\toprule
\textbf{Mô hình} & \textbf{R@1} & \textbf{R@5} & \textbf{R@10} & \textbf{\mrr} & $\alpha$ \\
\midrule
\multicolumn{6}{@{}l}{\textit{Visual Genome}} \\
CLIP thuần & 0,0584 & 0,1422 & 0,2018 & 0,1076 & 1 \\
Đối chứng & 0,0697\pmm0,0008 & \textbf{0,1742}\pmm0,0017 & \textbf{0,2485}\pmm0,0007 & 0,1277\pmm0,0001 & 0,1--0,2 \\
G-Retriever & \textbf{0,0708}\pmm0,0006 & 0,1724\pmm0,0011 & 0,2474\pmm0,0001 & \textbf{0,1278}\pmm0,0005 & 0,2--0,3 \\
\midrule
\multicolumn{6}{@{}l}{\textit{Flickr30k}} \\
CLIP thuần & 0,4901 & 0,7426 & 0,8250 & 0,6045 & 1 \\
Đối chứng cùng cấu hình & \textbf{0,5314}\pmm0,0019 & \textbf{0,7805}\pmm0,0011 & 0,8548\pmm0,0011 & \textbf{0,6430}\pmm0,0016 & 0,5 \\
G-Retriever & 0,5278\pmm0,0002 & 0,7794\pmm0,0006 & \textbf{0,8555}\pmm0,0004 & 0,6408\pmm0,0003 & 0,5 \\
\bottomrule
\end{tabular}}
\end{table}

\begin{table}[t]
\centering
\caption{G-Retriever giai đoạn 2 trên val (tập dùng để chọn epoch và $\alpha$), 3 seed; đối chứng như Bảng~\ref{tab:gret2}}
\label{tab:gret2-val}
\scriptsize
\setlength{\tabcolsep}{2pt}
\resizebox{\columnwidth}{!}{\begin{tabular}{@{}l c c c c@{}}
\toprule
\textbf{Mô hình} & \textbf{R@1} & \textbf{R@5} & \textbf{R@10} & \textbf{\mrr} \\
\midrule
\multicolumn{5}{@{}l}{\textit{Visual Genome}} \\
CLIP thuần & 0,0559 & 0,1443 & 0,2046 & 0,1061 \\
Đối chứng & 0,0680\pmm0,0000 & 0,1731\pmm0,0011 & 0,2406\pmm0,0005 & 0,1246\pmm0,0003 \\
G-Retriever & 0,0680\pmm0,0003 & 0,1731\pmm0,0008 & 0,2423\pmm0,0012 & 0,1244\pmm0,0003 \\
\midrule
\multicolumn{5}{@{}l}{\textit{Flickr30k}} \\
CLIP thuần & 0,4891 & 0,7368 & 0,8168 & 0,6016 \\
Đối chứng & 0,5281\pmm0,0017 & 0,7778\pmm0,0011 & 0,8513\pmm0,0010 & 0,6397\pmm0,0015 \\
G-Retriever & 0,5259\pmm0,0003 & 0,7768\pmm0,0002 & 0,8528\pmm0,0004 & 0,6381\pmm0,0002 \\
\bottomrule
\end{tabular}}
\end{table}

\begin{table}[t]
\centering
\caption{\mrr{} G-Retriever giai đoạn 2 trừ đối chứng cùng cấu hình, ghép cặp theo seed}
\label{tab:gret2-paired}
\scriptsize
\setlength{\tabcolsep}{2pt}
\begin{tabular}{@{}l c c c c@{}}
\toprule
\textbf{Tập} & \textbf{Seed 0} & \textbf{Seed 1} & \textbf{Seed 2} & \textbf{Mean \pmm{} std} \\
\midrule
VG, test & $-0,00042$ & +0,00031 & +0,00053 & +0,00014\pmm0,00050 \\
Flickr30k, test & $-0,00372$ & $-0,00119$ & $-0,00156$ & $-0,00216$\pmm0,00137 \\
Flickr30k, val & $-0,00340$ & $-0,00050$ & $-0,00115$ & $-0,00168$\pmm0,00152 \\
\bottomrule
\end{tabular}
\end{table}

Trên Visual Genome, giai đoạn 2 ngang đối chứng cùng cấu hình: hiệu \mrr{} test ghép cặp là +0,00014 \pmm{} 0,00050 và đổi dấu giữa các seed; trên val hiệu trung bình là $-0,0002$. Chỉ R@1 test cao hơn (+0,0011), nhưng R@5 và R@10 lại thấp hơn, nên không có chỉ số nào cho thấy đồ thị con giúp ổn định. Trên Flickr30k, cạnh đồng xuất hiện được dùng lại từ train Visual Genome (không dùng thống kê nào của Flickr30k), và giai đoạn 2 \emph{thấp hơn} đối chứng ở cả 3 seed, trên cả val và test (test: $-0,0022$ \pmm{} 0,0014). Cả hai mô hình vẫn vượt CLIP thuần rõ rệt, nên mức tăng một lần nữa đến từ phép chiếu học được. $\alpha$ chọn được gần với đối chứng (0,2--0,3 so với 0,1--0,2 trên Visual Genome; 0,5 như nhau trên Flickr30k), tức nhánh đồ thị con không làm mô hình bớt dựa vào CLIP. Như trong phần ablation, std ở đây chỉ đo dao động theo seed với 3 seed, không phải kiểm định theo truy vấn.
 Theo quy tắc dừng đã chốt trước, nhóm ghi nhận đây là kết quả âm: đồ thị con phụ thuộc truy vấn không mang thêm thông tin so với phép chiếu học được. Mã nguồn và bảng số của giai đoạn 2 nằm trong kho mã (\code{g\_retriever\_test.csv}, \code{g\_retriever\_flickr30k\_test.csv}, ADR-0051 và ADR-0056).

\subsection{Bộ truy vấn trực tiếp và gián tiếp}
\label{sec:kgbench}
Đề yêu cầu một bộ ít nhất 20 truy vấn có nhóm gián tiếp, kiểu \emph{child -- plays -- football}, để đánh giá riêng vai trò của đồ thị. Phiên bản đầu sinh truy vấn từ scene graph: câu gián tiếp giấu đối tượng ở giữa hai quan hệ, ví dụ ``A person is holding something that is in a sky''. Pilot trên val cho thấy kiểu câu này mơ hồ, vì ``something'' khớp với quá nhiều thứ. Phiên bản dùng trong báo cáo (ADR-0057) vì vậy giấu một \emph{khái niệm} và mô tả nó bằng tri thức ConceptNet~\cite{speer2017conceptnet}.

\textbf{Cách sinh.} Khái niệm đích là một hoạt động (\code{playing}, \code{riding}, \code{flying}) hoặc một đồ vật. Câu trực tiếp gọi tên khái niệm (``A person playing tennis.''). Câu gián tiếp chỉ dùng một từ chỉ loại và hai tính chất ConceptNet, sao cho trong kho khái niệm so sánh chỉ khái niệm đích thỏa cả ba (``A person playing a sport that has an advantage and uses a ball.''). Cạnh tính chất phải có trọng số ConceptNet $\ge 2$. Khái niệm đích cần ít nhất 20 ảnh train. Đáp án ứng viên là các ảnh test có annotation của khái niệm, từ 3 đến 50 ảnh. Kho so sánh và ngưỡng ảnh train chỉ đếm trên train. Mọi luật được chốt sau ba lần pilot trên val, trước khi sinh trên test; từ đó config không đổi giá trị nào.

\textbf{Duyệt bởi người.} Trên test, bộ sinh cho 28 cặp (8 hoạt động, 20 đồ vật). Một thành viên duyệt cả 28 cặp trên một trang có ảnh (ADR-0058). Trang chỉ hiện câu, khái niệm bị giấu, các cạnh ConceptNet và ảnh ứng viên, không hiện kết quả của mô hình nào. Với mỗi cặp, người duyệt giữ những ảnh thật sự khớp câu trực tiếp, hoặc bỏ cả cặp kèm lý do. Bảng~\ref{tab:kgbench} tóm tắt kết quả. Sáu cặp bị bỏ vì gần như không có ảnh ứng viên nào thật sự chứa khái niệm, ví dụ \code{measure}, \code{cart}, \code{jar}. Ba cặp bị bỏ vì câu gián tiếp mơ hồ hoặc lệch nghĩa, ví dụ ``pan'' bị hiểu thành cái khay. Trong 19 cặp được giữ, người duyệt giữ 291/426 ảnh ứng viên (68\%). Như vậy annotation Visual Genome thường gắn nhãn khái niệm cho cả ảnh mà khái niệm chỉ xuất hiện rất nhỏ hoặc không đúng nghĩa.

\begin{table}[t]
\centering
\caption{Bộ truy vấn KG trên test Visual Genome sau khi duyệt}
\label{tab:kgbench}
\scriptsize
\setlength{\tabcolsep}{3pt}
\begin{tabular}{@{}l c c c@{}}
\toprule
 & \textbf{Hoạt động} & \textbf{Đồ vật} & \textbf{Tổng} \\
\midrule
Cặp sinh trên test & 8 & 20 & 28 \\
Cặp bị bỏ & 1 & 8 & 9 \\
Cặp được giữ & 7 & 12 & \textbf{19} \\
Truy vấn (trực tiếp + gián tiếp) & 14 & 24 & \textbf{38} \\
\midrule
\multicolumn{4}{@{}l}{Ảnh đúng mỗi cặp: trung vị 12 (3--37); giữ 291/426 ảnh ứng viên} \\
\multicolumn{4}{@{}l}{Độ dài câu: trực tiếp 2,9 từ, gián tiếp 12,5 từ (trung bình)} \\
\bottomrule
\end{tabular}
\end{table}

\textbf{Giao thức đánh giá.} Bộ truy vấn chỉ dùng để đánh giá, không chọn hay chỉnh gì trên nó. Mỗi câu được mã hóa bằng đúng CLIP đóng băng và lấy 100 ảnh test gần nhất, như giao thức chính. Mỗi mô hình xếp lại 100 ảnh này bằng đúng checkpoint và $\alpha$ đã chọn trên val, không đọc khái niệm bị giấu hay các cạnh ConceptNet của câu. Nhóm chấm CLIP thuần, đối chứng không đồ thị cùng cấu hình (lr $10^{-4}$) và G-Retriever giai đoạn 2. Mỗi mô hình huấn luyện được chấm với 3 seed, báo cáo R@1/5/10 và \mrr{} tách riêng trực tiếp và gián tiếp, và tách theo nhóm hoạt động và đồ vật. Code (\code{scripts/evaluate\_query\_benchmark\_models.py}) đã có và đã qua kiểm tra rò rỉ. Bảng~\ref{tab:kgbench-res} cho kết quả. Chấm thêm các mô hình cấp độ 1 (CLIP + Graph, mô hình có cạnh đồng xuất hiện) trên bộ này là một hướng phát triển (Mục~\ref{sec:conclusion}).

\begin{table}[t]
\centering
\caption{\mrr{} trên bộ 38 truy vấn KG (test Visual Genome), mean \pmm{} std qua 3 seed; CLIP thuần tất định. Đối chứng: không đồ thị, cùng cấu hình huấn luyện (lr $10^{-4}$)}
\label{tab:kgbench-res}
\scriptsize
\setlength{\tabcolsep}{2pt}
\resizebox{\columnwidth}{!}{\begin{tabular}{@{}l c c c c c@{}}
\toprule
\textbf{Mô hình} & \textbf{Trực tiếp} & \textbf{Gián tiếp} & \textbf{GT -- hoạt động} & \textbf{GT -- đồ vật} & \textbf{Tất cả} \\
 & (19) & (19) & (7) & (12) & (38) \\
\midrule
CLIP thuần & 0,531 & 0,205 & 0,341 & 0,126 & 0,368 \\
Đối chứng & \textbf{0,585}\pmm0,037 & 0,331\pmm0,013 & \textbf{0,423}\pmm0,018 & 0,277\pmm0,027 & 0,458\pmm0,015 \\
G-Retriever GĐ2 & 0,559\pmm0,010 & \textbf{0,377}\pmm0,032 & 0,368\pmm0,048 & \textbf{0,382}\pmm0,028 & \textbf{0,468}\pmm0,021 \\
\bottomrule
\end{tabular}}
\end{table}

\textbf{Đọc kết quả.} Trên cả hai nhóm câu, cả hai mô hình có học đều vượt CLIP thuần. Mức tăng ở câu gián tiếp lớn hơn nhiều so với câu trực tiếp (0,205 lên 0,33--0,38, so với 0,531 lên 0,56--0,59). Phần lớn mức tăng này đã có ở đối chứng không đồ thị, tức là đến từ phép chiếu học được, giống kết luận trên giao thức chính. So với đối chứng, G-Retriever cao hơn ở câu gián tiếp nói chung (hiệu ghép cặp +0,046 \pmm{} 0,042, 2/3 seed), rõ nhất ở câu gián tiếp về đồ vật, cao hơn ở cả 3 seed (+0,105 \pmm{} 0,051), nhưng thấp hơn ở câu gián tiếp về hoạt động ($-0,055$ \pmm{} 0,061) và ở câu trực tiếp ($-0,026$ \pmm{} 0,028). Trên toàn bộ 38 câu, hai mô hình gần như ngang nhau (+0,010 \pmm{} 0,009). Như vậy chưa có bằng chứng ổn định rằng đồ thị con PCST giúp truy vấn gián tiếp nói chung. Riêng nhóm đồ vật gián tiếp có xu hướng được lợi, nhưng nhóm này chỉ có 12 câu nên cần thêm câu mới kết luận được.

\textbf{Giới hạn.} Đáp án do một người kiểm, không có người thứ hai đối chiếu. Đáp án ứng viên lấy từ annotation, nên ảnh đúng mà annotation bỏ sót sẽ không được tính, và recall là cận dưới. Bộ chỉ có 38 câu, nên chênh lệch giữa các mô hình trên bộ này chỉ là xu hướng. Riêng trên MMKG (Mục~\ref{sec:mmkg}), nhóm đã có bộ 318 truy vấn trực tiếp và gián tiếp được đánh giá đầy đủ.
